\section{Descent of nef and flat lines}\label{sec:descent}

A nef class that vanishes on the fibers of a fibration should come from
its base.  Singular fibers make this assertion more delicate than
cohomological descent on a smooth family: the discrepancy can contain
vertical divisors.  We first remove those divisors, and then distinguish
numerical descent from descent of an actual rational line.  Throughout
this section a \emph{fibration} is a surjective holomorphic map with
connected fibers.

\begin{proposition}[Numerical descent of a nef line]\label{desc:nef}
Let $f:X\to Y$ be a fibration of smooth connected compact K\"ahler
manifolds, and let $M\in\Pic(X)\otimes\Q$ be analytically nef.
Suppose that $c_1(M|_{X_y})=0$ on a general smooth fiber.
There are smooth compact K\"ahler modifications $p:X'\to X$ and
$q:Y'\to Y$, a fibration $g:X'\to Y'$ satisfying $fp=qg$, and a
rational line $N$ on $Y'$ such that
\begin{equation}\label{desc:numerical-identity}
                 p^*M\num g^*N.
\end{equation}
If $Y$ is projective, $Y'$ can be chosen projective and $N$ is nef.
\end{proposition}

\begin{proof}
We descend a positive current over the smooth fibers, remove the
vertical divisor discrepancy over the remaining fibers, and finally
recover the rationality of the descended class.
If $Y$ is a point, the hypothesis says that $c_1(M)=0$.
If the relative dimension is zero, $f$ is bimeromorphic; take a common
smooth model as both $X'$ and $Y'$.  We may therefore assume
$k=\dim Y>0$ and $d=\dim X-\dim Y>0$.

Flatten $f$ after a smooth modification $q:Y'\to Y$
\cite{HironakaFlattening}.
Let $Z$ be the normalization of the main component of $X\times_Y Y'$,
and resolve it by a projective modification $r:X'\to Z$.
Write $g_0:Z\to Y'$ and $g=g_0r$.  The map $g_0$ is equidimensional;
it and $g$ have connected fibers by Stein factorization and normality
of $Y'$.  All these spaces admit the asserted K\"ahler models.
Write $p_0:Z\to X$ for the map to the original source, so that
$p=p_0r$ is a modification.  If $Y$ is projective, the base
modifications may be chosen projective.  Figure~\ref{desc:diagram}
separates the equidimensional map from the resolution above it.

\begin{figure}[ht]
\centering
\begin{tikzpicture}[>=Stealth, every node/.style={inner sep=3pt}]
\node (Xp) at (0,1.35) {$X'$};
\node (Z) at (2.6,1.35) {$Z$};
\node (X) at (5.2,1.35) {$X$};
\node (Yp) at (2.6,0) {$Y'$};
\node (Y) at (5.2,0) {$Y$};
\draw[->] (Xp) -- node[above] {$r$} (Z);
\draw[->] (Z) -- node[above] {$p_0$} (X);
\draw[->] (Xp) -- node[below left] {$g$} (Yp);
\draw[->] (Z) -- node[right] {$g_0$} (Yp);
\draw[->] (X) -- node[right] {$f$} (Y);
\draw[->] (Yp) -- node[below] {$q$} (Y);
\end{tikzpicture}
\caption{The flattened main component $Z$ is equidimensional over
$Y'$.  A prime divisor on the resolution $X'$ whose image in $Y'$
has codimension at least two must therefore be exceptional for $r$.}
\label{desc:diagram}
\end{figure}

Put $\alpha=c_1(p^*M)$, choose a K\"ahler form $\omega$ on $X'$, and
choose a closed positive current $T$ representing $\alpha$.
The volume
\[
                  c=\int_{X'_y}\omega^d
\]
is positive and constant on the smooth-fibration locus.  Define the
closed positive $(1,1)$-current
\begin{equation}\label{desc:base-current}
                  S=c^{-1}g_*(T\wedge\omega^d)
\end{equation}
on $Y'$.  We claim that $T=g^*S$ over that locus.
Indeed $g_*(T\wedge\omega^{d-1})$ is a closed positive current of
degree zero, hence a nonnegative constant.  Its cohomology class is
the fiber intersection of $\alpha$ with $\omega^{d-1}$, which vanishes.
Thus this current is zero.  In local product coordinates this says
that the vertical block of the positive matrix of measures defining
$T$ vanishes.  Positivity also forces its mixed block to vanish.
Closedness then makes the horizontal coefficients independent of the
fiber coordinates.  Connectedness of the fibers makes these local
currents descend, and Equation~\eqref{desc:base-current} identifies
their common descent as $S$.

The current $S$ has local plurisubharmonic potentials, so its pullback
by the dominant map $g$ is defined on all of $X'$.
The closed order-zero current $T-g^*S$ is supported on the inverse
image of the complement of the smooth-fibration locus in $Y'$.
The support theorem for currents
\cite[Chapter~III, Corollaries~2.11 and~2.14]{DemaillyCADG} therefore gives
\begin{equation}\label{desc:vertical-difference}
              T-g^*S=[D],\qquad D=\sum_E x_E E,
\end{equation}
where $D$ is a signed real divisor supported on vertical primes.
We next show that its components dominating base divisors come in
whole pullbacks.

Fix a prime divisor $Q\subset Y'$ occurring among their images, and
let $E_1,\ldots,E_s$ be the primes of $X'$ dominating $Q$.
Let $a_i>0$ be their multiplicities in $g^*Q$, and let $x_i$ be their
coefficients in $D$.  For a K\"ahler form $\eta$ on $Y'$, set
\begin{equation}\label{desc:intersection-matrix}
       C_{ij}=\int_{X'}\{E_i\}\{E_j\}\,(g^*\eta)^{k-1}\omega^{d-1}.
\end{equation}
For $i\ne j$, these numbers are nonnegative.
The inverse image $g^{-1}(Q)$ is an effective Cartier divisor, hence
has pure codimension one.  After a generic choice of the point in
$Q$, its fiber components all have pure dimension $d$.  Distinct
primes $E_i,E_j$ have pure codimension-two intersection; if they
meet over a general point of $Q$, that intersection dominates $Q$
and has fibers of dimension $d-1$.  It therefore contributes
positively to $C_{ij}$.  These dimension statements follow by
generic flatness over $Q$, excluding the smaller base images.
The graph having an edge when $C_{ij}>0$ is consequently connected.  Indeed, over a general point
of $Q$ the irreducible components of the connected fiber have a
connected intersection graph.  Grouping those components according
to the global prime $E_i$ containing them gives the graph above as
a quotient, which is still connected.  Moreover $Ca=0$, where
$a=(a_i)$.
To see this, pair $E_i$ with $g^*Q$ and $(g^*\eta)^{k-1}\omega^{d-1}$.
Its image lies in $Q$, so the $k$ classes from the base give zero.
Components of $g^*Q$ whose images have codimension at least two give
zero separately in this pairing.

Let $\beta$ be the Bott--Chern class of $S$.  Equation~\eqref{desc:vertical-difference} gives
$\{D\}=\alpha-g^*\beta$.  Pairing this identity with each
$E_i(g^*\eta)^{k-1}\omega^{d-1}$ shows that
\[
                               Cx\geq0.
\]
Here the $\alpha$ term is nonnegative by nefness, and the base term
vanishes by dimension.  Components of $D$ above other base divisors,
or above sets of codimension at least two, also contribute zero.
Since $a$ has strictly positive entries and $a^{\mathsf t}Cx=0$,
we obtain $Cx=0$.  The kernel is precisely $\R a$: indeed
\[
 x^{\mathsf t}Cx
 =-\sum_{i<j} C_{ij}a_i a_j
       \left(\frac{x_i}{a_i}-\frac{x_j}{a_j}\right)^2,
\]
and the graph is connected.  Consequently $x_i=t_Qa_i$ for one
real number $t_Q$.

Subtract $g^*(\sum_Q t_QQ)$ from $D$, and denote the resulting divisor
by $D_0$.  Its image in $Y'$ has codimension at least two.
Equidimensionality of $g_0$ implies that every such prime on $X'$ is
$r$-exceptional.  Also
\[
 \{D_0\}=\alpha-g^*\beta',\qquad
 \beta'=\beta+\sum_Q t_Q\{Q\}.
\]
Both classes on the right come from classes with local potentials
on $Z$: $\alpha$ comes from the original source $X$, and $\beta'$
comes from $Y'$.  Thus $D_0$ has degree zero on every curve contracted
by $r$.  Relative negativity for a projective morphism of normal analytic
spaces \cite[Section~11, after Definition~11.1]{FujinoAnalyticBCHM}, applied to the exceptional divisors
$D_0$ and $-D_0$, gives $D_0=0$.  We have proved
\begin{equation}\label{desc:real-pullback}
                              \alpha=g^*\beta'.
\end{equation}

It remains to justify that the descended class is a rational line
class; the use of $\omega$ above does not supply rationality by
itself.  Pullback
\[
                    g^*:H^2(Y',\Q)\longrightarrow H^2(X',\Q)
\]
is injective.  Over $\R$ a left inverse is obtained by pushing
against $\omega^d$ and dividing by $c$, so injectivity follows there
and hence over $\Q$.  An injective rational linear map has a rational
preimage for every rational vector in its real image.  Since
$\alpha$ is rational, Equation~\eqref{desc:real-pullback} therefore
makes $\beta'$ rational.  It has type $(1,1)$, and the Lefschetz
$(1,1)$ theorem gives a rational line $N$ with $c_1(N)=\beta'$.

Finally suppose that $Y'$ is projective.  For any irreducible curve
$C\subset Y'$, resolve a component of $g^{-1}(C)$ that dominates
$C$, and factor its map through the normalization of $C$.
The pullback of $\alpha$ is nef.  Pairing it with a K\"ahler power
on that resolution gives a positive constant times $\deg(N|_C)$.
This degree is nonnegative.  The projective numerical criterion for
nefness now shows that $N$ is nef.
\end{proof}

The preceding proposition descends the Chern class.  For later use we
also need an actual rational-line identity.  The price is a single
flat twist on the original source.

\begin{corollary}[Choosing a descended representative]
\label{desc:representative}
In the notation and under the hypotheses of
Proposition~\ref{desc:nef}, there is a rational line $M^*$ on $X$
with $M^*\num M$ such that
\[
                          p^*M^*\qsim g^*N.
\]
\end{corollary}

\begin{proof}
The rational line $g^*N-p^*M$ has zero real Chern class.  After clearing
denominators and torsion, it belongs to $\Pic^0(X')$.
Lemma~\ref{pre:picard}(2) identifies $\Pic^0(X)$ with
$\Pic^0(X')$.  Thus
$g^*N-p^*M\qsim p^*F$ for some $F\in\Pic^0(X)\otimes\Q$.
Take $M^*=M+F$.
\end{proof}

The next lemma explains which fiberwise trivial flat twists already
come from the base.  Passing to rational lines allows us to remove
finite meridian holonomy and torsion Chern classes.

\begin{lemma}[Descent of a flat line]\label{desc:flat}
Let $f:X\to Y$ be a fibration of smooth connected compact K\"ahler
manifolds, and let $F\in\Pic^0(X)\otimes\Q$.
If $F$ is trivial as a rational line on a general smooth fiber, then
\[
                         F\qsim f^*G
\]
for some $G\in\Pic^0(Y)\otimes\Q$.
\end{lemma}

\begin{proof}
Take a positive multiple so that $F$ is a genuine unitary flat line
and its restriction to one smooth fiber $X_y$ is holomorphically
trivial.  On a compact connected K\"ahler manifold, a holomorphically
trivial unitary flat line has trivial holonomy.  Thus the character
of $F$ is trivial on $\pi_1(X_y)$.

Choose a dense Zariski open $Y^\circ$ over which $f$ is smooth, and
write $X^\circ=f^{-1}(Y^\circ)$.  The homotopy sequence of the smooth
proper fibration gives
\[
 \pi_1(X_y)\longrightarrow\pi_1(X^\circ)
       \longrightarrow\pi_1(Y^\circ)\longrightarrow1.
\]
The character of $F|_{X^\circ}$ therefore comes from a unitary
character $\chi$ of $\pi_1(Y^\circ)$.

Let $Q$ be a divisorial component of $Y\setminus Y^\circ$.
Choose a prime dominating $Q$ and a transverse disk at a general
point of that prime, away from the other components of the inverse
image of $Y\setminus Y^\circ$.  Its image winds $a_Q$ times around a small meridian
of $Q$, where $a_Q>0$ is the multiplicity of that prime in $f^*Q$.
The disk lies in $X$, so the character of $F$ is trivial on its
boundary.  Hence $\chi$ has finite order dividing $a_Q$ on the base
meridian.  There are finitely many such $Q$.  A common power of
$\chi$ kills all their meridians and therefore factors through
$\pi_1(Y)$; subsets of complex codimension at least two add no
obstruction.  The resulting unitary flat line on $Y$ has zero real
Chern class, and another power removes any torsion in its integral
Chern class.  It then belongs to $\Pic^0(Y)$.

The resulting characters agree after pullback on $X^\circ$.
Since $\pi_1(X^\circ)\to\pi_1(X)$ is surjective, they agree on $X$
as well.  Dividing by the powers taken in the argument proves the
asserted identity in $\Pic(X)\otimes\Q$.
\end{proof}
